\chapter{Metodología}

El proyecto corresponde a una \textbf{investigación aplicada, cuantitativa,
experimental y de desarrollo tecnológico}. Su propósito es reactivar el robot
cuadrúpedo Nova Spot Micro y evaluar, mediante mediciones repetibles, el efecto
de una política de aprendizaje por refuerzo sobre una estrategia convencional
de locomoción.

La arquitectura de control tendrá dos niveles. El primero generará una marcha
nominal mediante una máquina de estados finitos, trayectorias de los pies y
referencias articulares para postura, paso y gateo. El segundo estará compuesto
por una política de aprendizaje por refuerzo que producirá únicamente
correcciones pequeñas y acotadas sobre dichas referencias. La política no
controlará directamente las señales PWM de los servomotores y permanecerá
subordinada a un supervisor independiente de seguridad, límites articulares y
saturaciones.

El desarrollo seguirá un esquema secuencial con retroalimentación. Cada etapa
deberá superar sus verificaciones antes de avanzar a pruebas que impliquen un
mayor riesgo para el hardware.

\section{Fases del desarrollo}

La metodología se organiza en ocho fases principales.

\subsection{Fase 1: Caracterización de la plataforma}

Se realizará el inventario técnico y la inspección mecánica y eléctrica del
robot. Se medirán la geometría, los límites articulares, los ceros mecánicos,
las masas y la ubicación aproximada del centro de masa. Las dimensiones
relevantes se medirán al menos tres veces y se reportarán mediante promedio e
incertidumbre. También se identificarán las variables que pueden medirse en el
robot real y las limitaciones de los actuadores y sensores disponibles.

\textbf{Resultado esperado:} ficha técnica, tabla de parámetros geométricos e
inerciales, inventario de componentes y lista de riesgos.

\subsection{Fase 2: Modelado y simulación}

Se formularán la cinemática directa e inversa, el jacobiano, los límites
articulares y el espacio de trabajo de cada pata. El modelo se extenderá a las
cuatro extremidades y se validará comparando posiciones calculadas y medidas.
Posteriormente se construirá el modelo del robot en URDF y se integrará con ROS
2 y MuJoCo, incluyendo masas, inercias, colisiones, fricción, actuadores e IMU
simulada.

\textbf{Resultado esperado:} modelo matemático validado y simulación
reproducible capaz de permanecer de pie y recibir referencias articulares.

\subsection{Fase 3: Arquitectura electrónica y seguridad}

Se dimensionarán la fuente de alimentación, la distribución de potencia, las
protecciones y el cableado a partir de la corriente máxima prevista. Se definirá
la comunicación entre el computador embebido, el controlador PWM, la IMU y los
demás sensores. Los servomotores utilizarán una alimentación de potencia
separada de la electrónica de control, con tierra común, fusible, interruptor y
parada de emergencia. Cada articulación se calibrará individualmente antes de
integrar las patas.

\textbf{Resultado esperado:} diagramas eléctricos, presupuesto energético,
archivo de calibración de las articulaciones y protocolos de encendido, apagado
y emergencia.

\subsection{Fase 4: Control convencional y marcha nominal}

Se implementarán el control de postura, la interpolación suave de referencias,
las trayectorias de los pies y una máquina de estados finitos para coordinar las
fases de apoyo y oscilación. Esta línea base ejecutará postura, paso y gateo en
una superficie plana. Incluirá límites de posición y velocidad, detección de
referencias no alcanzables y un supervisor de seguridad.

No se fija una meta de velocidad máxima. El robot operará a baja velocidad y
esta se limitará según la estabilidad observada, la capacidad de los actuadores
y la seguridad. La velocidad obtenida se medirá y reportará, mientras que el
criterio funcional será completar ciclos de locomoción verificables y
repetibles.

\textbf{Resultado esperado:} robot simulado ejecutando postura, paso y gateo
sin aprendizaje por refuerzo, como condición de referencia para la comparación.

\subsection{Fase 5: Formulación y entrenamiento del aprendizaje por refuerzo}

El problema se formulará como un proceso de decisión de Markov. Las
observaciones incluirán, según la disponibilidad de sensado, orientación y
velocidad angular de la IMU, posiciones articulares, fase de marcha,
referencias nominales y contactos estimados o medidos. Las acciones serán
correcciones limitadas de las referencias articulares o de las posiciones de
los pies. La recompensa considerará avance, estabilidad, seguimiento,
suavidad, esfuerzo y penalización por caída o saturación.

Se evaluará inicialmente el algoritmo PPO por su compatibilidad con espacios de
acción continuos. El entrenamiento se realizará primero en MuJoCo, empleando
aleatorización de masa, centro de masa, fricción, ganancias, retardos, ruido de
sensores y capacidad de los actuadores. Se conservarán semillas, curvas de
recompensa, hiperparámetros, configuraciones y políticas entrenadas.

\textbf{Resultado esperado:} formulación matemática completa, protocolo de
entrenamiento reproducible y política correctiva validada en simulación.

\subsection{Fase 6: Integración y transferencia al robot}

La transferencia se realizará progresivamente: ejecución sin activar los
servos, servo individual, pata suspendida, cuatro patas con el robot suspendido,
postura con soporte, transferencia de peso, un paso, ciclo completo de gateo y
marcha continua. Solo después se habilitarán las correcciones aprendidas, con
amplitud limitada. En todas las etapas se mantendrán la parada de emergencia y
los límites independientes de la política.

\textbf{Resultado esperado:} ejecución segura de la marcha nominal y de la
marcha con corrección aprendida en el robot físico.

\subsection{Fase 7: Validación experimental comparativa}

Se evaluarán dos condiciones principales: marcha nominal sin aprendizaje y la
misma marcha con la capa correctiva aprendida. Ambas se probarán en simulación y
en el robot físico sobre una superficie plana y en condiciones controladas. Se
realizarán al menos diez ciclos válidos por modo de marcha y condición.

Las variables de respuesta serán el margen de estabilidad estática, la
inclinación del cuerpo, el desplazamiento, la desviación lateral, la longitud y
el tiempo de paso, la repetibilidad ciclo a ciclo y el número de fallos o
caídas. El seguimiento articular se evaluará mediante RMSE, error máximo,
retardo, sobreimpulso y error en estado estable únicamente si existe medición
real de la posición; el controlador PWM por sí solo no proporciona esta
realimentación.

\textbf{Resultado esperado:} conjunto de datos y comparación estadística que
permita determinar en qué medida la política mejora la estabilidad y la
repetibilidad respecto de la marcha nominal.

\subsection{Fase 8: Análisis y documentación}

Se contrastarán los resultados de simulación y del robot físico, se analizarán
las limitaciones impuestas por el hardware y se establecerá la trazabilidad
entre cada objetivo, método, evidencia y conclusión. Se documentarán el código,
las configuraciones, el protocolo experimental y el manual de operación.

\textbf{Resultado esperado:} informe final reproducible, conclusiones
vinculadas con los objetivos y documentación técnica para reutilizar la
plataforma.

\section{Diagrama de flujo de la metodología}

\begin{figure}[h!]
\centering
\resizebox{0.96\textwidth}{!}{%
\begin{tikzpicture}[
    node distance=0.65cm,
    block/.style={rectangle, draw, fill=blue!20, text width=10em,
        text centered, rounded corners, minimum height=3.5em},
    arrow/.style={-latex, thick}
]
\node (f1) [block] {1. Caracterización de la plataforma};
\node (f2) [block, right=of f1] {2. Modelado y simulación};
\node (f3) [block, right=of f2] {3. Electrónica y seguridad};
\node (f4) [block, right=of f3] {4. Marcha nominal};
\node (f5) [block, below=of f4] {5. Aprendizaje por refuerzo};
\node (f6) [block, left=of f5] {6. Transferencia al robot};
\node (f7) [block, left=of f6] {7. Validación comparativa};
\node (f8) [block, left=of f7] {8. Análisis y documentación};

\draw[arrow] (f1) -- (f2);
\draw[arrow] (f2) -- (f3);
\draw[arrow] (f3) -- (f4);
\draw[arrow] (f4) -- (f5);
\draw[arrow] (f5) -- (f6);
\draw[arrow] (f6) -- (f7);
\draw[arrow] (f7) -- (f8);
\draw[-latex, dashed] (f7.north) .. controls +(up:1.2cm) and +(down:1.2cm) .. (f4.south);
\draw[-latex, dashed] (f6.north) .. controls +(up:1.0cm) and +(down:1.0cm) .. (f3.south);
\end{tikzpicture}%
}
\caption{Fases de desarrollo y ciclos de retroalimentación del proyecto.}
\label{fig:metodologia_fases}
\end{figure}

\section{Articulación con los objetivos}

La caracterización y el modelado responden al primer objetivo específico; la
arquitectura electrónica y de seguridad, al segundo; la marcha nominal, al
tercero; el entrenamiento y la transferencia de la política correctiva, al
cuarto; y la validación comparativa con y sin aprendizaje por refuerzo, al
quinto. Esta correspondencia permite relacionar cada conclusión con evidencia
experimental concreta.
